\documentclass[12pt]{article}
\usepackage[utf8]{inputenc}
\usepackage[T1]{fontenc}
\usepackage{hyperref}
\title{Marenje Language Resources for K-12 Curriculum Development}
\author{Compiled by Claude Sonnet 5 (claude-sonnet-5) for Andr\'as Kornai}
\date{}
\begin{document}
\maketitle
\section{Introduction}
Marenje (ISO 639-3 vmr, Glottocode mare1260) is a Bantu language of northern Mozambique near the Malawi border, with Milange and Tengua as its heartland. Glottolog places it in North Mozambique Bantu (Chuwaboic), whereas Ethnologue groups it under Makhuwa. ASJP lists 588,000 speakers, and Joshua Project reports a people group of 813,000. Joshua Project also says younger people increasingly prefer Chichewa or Portuguese, and that the language has no literature and no Bible.
\section{Primary Educational Resources}
None found.
\section{Linguistic and Reference Materials}
\begin{itemize}
\item Glottolog reference record for Alfazema, Cantau\'ile, Rijone and Kr\"oger (2006), a grammar sketch: \url{https://glottolog.org/resource/reference/id/94422}
\item Glottolog language page for Marenje: \url{https://glottolog.org/resource/languoid/id/mare1260.bigmap.html}
\item Glottolog reference record: \url{https://glottolog.org/resource/reference/id/577247}
\item ASJP wordlist: \url{https://asjp.clld.org/languages/MARENJE}
\end{itemize}
\section{Religious and Community Resources}
\begin{itemize}
\item Joshua Project people group page (Marenje in Mozambique): \url{https://www.joshuaproject.net/people_groups/19026/MZ}
\item Joshua Project language page: \url{https://www.joshuaproject.net/index.php/languages/vmr}
\item Scripture Earth listing, which per Joshua Project has no Marenje scripture: \url{https://www.scriptureearth.org/00eng.php?iso=vmr}
\end{itemize}
\section{Digital and Audio Resources}
None found.
\section{Curriculum Development Context}
The search found almost no educational materials. Churches use Chichewa or Portuguese. A mother-tongue curriculum would need new orthography, primers and teacher training, building on the 2006 grammar sketch.
\end{document}